\section{Background}
\label{sec:background}

\paragraph{Fault localisation.}
Automated fault localisation (FL) aims to rank program elements---statements,
methods, or functions---by their likelihood of containing a defect.
Spectrum-based FL (SBFL) methods such as Tarantula~\cite{jones05tarantula}
compute suspiciousness scores from pass/fail test-coverage data.
A broad survey of FL techniques is given by \citet{wong16survey}.
Our work differs from SBFL: we ask an LLM to identify the faulty function
from a natural-language bug description and optional call-graph context,
without access to test suites or execution traces.

\paragraph{LLM agents with tool use.}
The \textsc{ReAct} framework~\cite{yao23react} interleaves reasoning steps
with calls to external tools and observations of their results.
OpenAI-compatible APIs expose this pattern via a \emph{function-calling}
interface: the model selects from a schema-defined tool set and the host
dispatches the call.
We use this interface to let the LLM navigate a call-graph database without
loading the full graph into the prompt.

\paragraph{Call graphs.}
A call graph records direct-call edges between functions in a program.
A \emph{compile-time} (CT) call graph is extracted by parsing source with a
compiler front-end, capturing precise call sites with source-level provenance.
A \emph{syntactic} (SYN) call graph is extracted by pattern-matching on the
parse tree without full compilation, trading precision for ease of construction.
Neither representation models function-pointer calls, creating structural gaps
when control flow crosses a pointer boundary.
